\scanpage{107}
où on pose
\[
\mathbb{U}=\{z\in\mathbb{C}\mid |z|=1\}
\]
--- l'orientation de ce $\mathbb{U}$-torseur \inserted{(déduite de celle de $\mathbb{U}$)} est opposée de celle induite par $X_1$. \deleted{\unsure{$X$}} La surface holomorphe $\unsure{X|\widetilde X}$ \inserted{avec les \unsure{pts} de $C$ \unsure{marqués}} se récupère à partir de la surface holomorphe à bord $X_1$, et de la structure précédente de $\mathbb{U}$-torseur sur chacune des comp. connexes \inserted{$T_{\widetilde x}=\partial_{\widetilde x}X$} du bord $\partial X_1$, en notant que le cône sur $\partial_{\widetilde x}X$ \unsure{e|c} devient, grâce à cette structure, un disque \inserted{à bord} holomorphe, et par chirurgie holomorphe on peut le recoller à $X_1$, pour obtenir $X$, avec ses marques par les centres des disques respectifs.

Considérons, pour $x\in C$, les \deleted{\unsure{ensembles}} \inserted{deux pts} \deleted{des} $x'$ et $x''$ de $\widetilde C$ au dessus de $x$, et l'espace
\[
\operatorname{Isom}_{\mathbb{U}}^{-}(T_{x'},T_{x''})\simeq\operatorname{Isom}_{\mathbb{U}}(T_{x'}^{-1},T_{x''})
\]
des antiisomorphismes de $T_{x'}$ avec $T_{x''}$ (i.e. hom de torseurs compatibles avec l'\deleted{isomorphisme}\inserted{\unsure{anti-involution|antiinvolution}} $z\longmapsto z^{-1}:\mathbb{U}\simeq\mathbb{U}$), qui s'identifie aussi, plus symétriquement, au torseur $T_{x'}\wedge T_{x''}$, ou encore au $\mathbb{U}$-torseur
\[
T_x=\deleted{\unsure{\prod}}\mathop{\bigwedge}_{\widetilde x\in C_x}T_{\widetilde x}
\]
\margin{\uline{NB} $T_{\widetilde x}$ s'identifie aussi au \unsure{cercle tangent} à $\widetilde X$ en $\widetilde x$, i.e. à $T_{\widetilde X,\widetilde x}^{*}/\mathbb{R}^{*+}$ (ce qui est indép. du choix des $D_{\widetilde x}$)}

Considérons enfin le produit
\[
P_C=P_{C,X}=\prod_{x\in C}T_x
\]
qui est un torseur sous $\mathbb{U}^C$. \deleted{On} Pour \inserted{$\mathcal{X}=$}$(X,(D_{\widetilde x}))$

\mspage{106}
